% Unit 049: English translation of chapter4.tex, lines 62--208.
% Source: Methods of Algebra, Volume 2, commit
% 9a5803ff77dd3257484cb177f851a73770a59dd3 (CC BY 4.0).
% stable-unit-id: o014.aljabr2.chapter4.triangulated-category-definition
% Coverage: Section 4.1, from categories with translation through the remark
% on duality; ends before the lemma at line 210.

% segment-id: o014.li2.chapter4.triangulated-category-definition.h001
\section{Definition of a Triangulated Category}\label{sec:triangular-def}
\hypertarget{o014.aljabr2.chapter4.triangulated-category-definition}{ }

% segment-id: o014.li2.chapter4.triangulated-category-definition.p001
A triangulated category is an additive category with translation, equipped
with a class of distinguished triangles. We begin by explaining what a
category with translation is.

% segment-id: o014.li2.chapter4.triangulated-category-definition.d001
\begin{definition}\label{def:category-with-translation}
	\index{category!with translation}
	\index{translation functor}
	A \emph{category with translation} is data $(\mathcal{D},T)$ consisting
	of a category $\mathcal{D}$ and an equivalence
	$T:\mathcal{D}\to\mathcal{D}$.%
	\footnote{Many references require $T$ to be an automorphism. In that case
	one may choose an inverse functor $T^{-1}$ satisfying
	$T^{-1}T=\identity_{\mathcal D}=TT^{-1}$. This condition covers every
	situation treated in this book and avoids possible subtleties from
	$2$-category theory, so the reader may harmlessly make the same
	assumption. The study of stable model categories or stable
	$\infty$-categories, however, cannot avoid the case in which the
	equivalence is not an isomorphism, especially in topology.}
	The equivalence $T$ is called the \emph{translation functor} of
	$\mathcal{D}$. In addition, choose a quasi-inverse $T^{-1}$ of $T$ and
	isomorphisms $T^{-1}T\simeq\identity_{\mathcal D}\simeq TT^{-1}$ so that
	these data form an adjoint equivalence in the sense of
	\cite[Theorem~2.6.12]{Li1}.

% segment-id: o014.li2.chapter4.triangulated-category-definition.l001
	\begin{itemize}
% segment-id: o014.li2.chapter4.triangulated-category-definition.i001
		\item A functor from $(\mathcal{D},T)$ to
		$(\mathcal{D}',T')$ is a functor $F:\mathcal{D}\to\mathcal{D}'$
		together with a chosen isomorphism $FT\simeq T'F$. By convention, the
		isomorphism is suppressed and these data are abbreviated as
		$F:(\mathcal{D},T)\to(\mathcal{D}',T')$.
% segment-id: o014.li2.chapter4.triangulated-category-definition.i002
		\item A morphism between two functors
% segment-id: o014.li2.chapter4.triangulated-category-definition.g001
		$\begin{tikzcd}
			(\mathcal{D},T) \arrow[r,shift left,"F"]
			\arrow[r,shift right,"{F'}"'] & (\mathcal{D}',T')
		\end{tikzcd}$
		is a morphism $\varphi:F\to F'$ that makes the following diagram
		commute:
% segment-id: o014.li2.chapter4.triangulated-category-definition.g002
		\[\begin{tikzcd}
			FT \arrow[r,"\varphi T"] \arrow[d,"\sim"' sloped] &
			F'T \arrow[d,"\sim" sloped] \\
			T'F \arrow[r,"{T'\varphi}"'] & T'F'.
		\end{tikzcd}\]
% segment-id: o014.li2.chapter4.triangulated-category-definition.i003
		\item A subcategory of a category with translation
		$(\mathcal{D},T)$ is a subcategory $\mathcal{D}'$ of $\mathcal{D}$
		that is closed under $T$ and for which
		$T':=T|_{\mathcal{D}'}$ makes $(\mathcal{D}',T')$ a category with
		translation.
	\end{itemize}

% segment-id: o014.li2.chapter4.triangulated-category-definition.p002
	If $\mathcal{D}$ is additive and $T$ is an additive functor (as it is
	automatically; see Corollary~\ref{prop:automatic-additivity}), then
	$(\mathcal{D},T)$ is called an \emph{additive category with translation};
	functors between such categories are also required to be additive. After
	fixing a commutative ring $\Bbbk$, the same definition applies in the
	$\Bbbk$-linear setting.
\end{definition}

% segment-id: o014.li2.chapter4.triangulated-category-definition.p003
If $(\mathcal{D},T)$ is a category with translation, then so is
$(\mathcal{D}^{\opp},T^{-1})$. By convention, the data
$(\mathcal{D},T)$ are often abbreviated to $\mathcal{D}$. The following
notion is particularly useful.

% segment-id: o014.li2.chapter4.triangulated-category-definition.d002
\begin{definition}[Morphisms with degree]\label{def:morphism-with-degree}
	\index[sym1]{mdegree@$X \xrightarrow{+m} Y$}
	A morphism of degree $m$ in a category with translation
	$(\mathcal{D},T)$ is a morphism of the form $f:X\to T^mY$, also written
	$f:X\xrightarrow{+m}Y$. Given
	$X\xrightarrow[f]{+m}Y\xrightarrow[g]{+n}Z$, their composite $gf$ is
	defined to be the morphism of degree $m+n$ given by $T^m(g)f$.
	Commutative diagrams can likewise be considered for morphisms of this
	kind.

	In the multivariable setting, suppose that $\mathcal{D}$ is equipped with
	strictly commuting automorphisms $T_1,\ldots,T_n$, so that
	$T_iT_j=T_jT_i$. A morphism of degree $(m_1,\ldots,m_n)$ is defined to be
	a morphism of the form $X\to T_1^{m_1}\cdots T_n^{m_n}Y$. It is written
	$X\xrightarrow{+(m_1,\ldots,m_n)}Y$.
\end{definition}

% segment-id: o014.li2.chapter4.triangulated-category-definition.p004
For example, given a category $\mathcal{A}$, the category
$(\mathcal{A}^{\Z},T)$ of graded objects from
Definition~\ref{def:graded-obj} is a category with translation. More
generally, one can consider the category $\mathcal{A}^{\Z^n}$ of
$\Z^n$-graded objects, together with a translation functor
$T_1,\ldots,T_n$ in each variable, and then discuss morphisms with degree
between $\Z^n$-graded objects. The multivariable case, however, is not a
focus of this chapter.
\index{graded object}

% segment-id: o014.li2.chapter4.triangulated-category-definition.ex001
\begin{example}\label{eg:cplx-as-translation-cat}
	The category of complexes $\cate{C}(\mathcal{A})$ over an additive
	category $\mathcal{A}$, together with the translation functor
	$T:X\mapsto X[1]$ from Definition~\ref{def:translation-functor}, is an
	additive category with translation. The category
	$\cate{K}(\mathcal{A})$ from Definition~\ref{def:cplx-homotopy-cat}
	inherits the translation functor from $\cate{C}(\mathcal{A})$, which is
	still denoted by $T$. The canonical functor
	$F:\cate{C}(\mathcal{A})\to\cate{K}(\mathcal{A})$ plainly satisfies
	$FT=TF$. The same applies, for every $\star\in\{+,-,\bdd\}$, to the
	categories $\cate{C}^{\star}(\mathcal{A})$ (or
	$\cate{K}^{\star}(\mathcal{A})$) introduced in
	Definition~\ref{def:cplx-cat-variant} (or
	Definition~\ref{def:cplx-homotopy-cat-variant}).
\end{example}

% segment-id: o014.li2.chapter4.triangulated-category-definition.p005
In this chapter we mainly consider additive categories with translation.

% segment-id: o014.li2.chapter4.triangulated-category-definition.d003
\begin{definition}
	\index{triangle}
	Let $(\mathcal{D},T)$ be an additive category with translation. A
	\emph{triangle} in it is a sequence of morphisms in $\mathcal{D}$ of the
	form
% segment-id: o014.li2.chapter4.triangulated-category-definition.q001
	\[ X\xrightarrow{f}Y\xrightarrow{g}Z\xrightarrow{h}TX. \]
	A morphism between two triangles is a commutative diagram of the form
% segment-id: o014.li2.chapter4.triangulated-category-definition.g003
	\[\begin{tikzcd}
		X \arrow[r,"f"] \arrow[d,"\alpha"'] &
		Y \arrow[r,"g"] \arrow[d,"\beta"] &
		Z \arrow[r,"h"] \arrow[d,"\gamma"] &
		TX \arrow[d,"T\alpha"] \\
		X' \arrow[r,"{f'}"'] & Y' \arrow[r,"{g'}"'] &
		Z' \arrow[r,"{h'}"'] & TX'.
	\end{tikzcd}\]
	The two rows are the given triangles. Such a morphism is also abbreviated
	to $(X,Y,Z)\to(X',Y',Z')$. Isomorphisms between triangles are defined in
	the same way.

% segment-id: o014.li2.chapter4.triangulated-category-definition.p006
	The same notion is available if the additive category is replaced by a
	more general $\Bbbk$-linear category, where $\Bbbk$ is any commutative
	ring.
\end{definition}

% segment-id: o014.li2.chapter4.triangulated-category-definition.p007
Clearly, triangles in $(\mathcal{D},T)$ are the same as triangles in
$(\mathcal{D}^{\opp},T^{-1})$.

% segment-id: o014.li2.chapter4.triangulated-category-definition.r001
\begin{remark}\label{rem:triangle-signs}
	Consider a triangle
	$X\xrightarrow{f}Y\xrightarrow{g}Z\xrightarrow{h}TX$ in a
	$\Bbbk$-linear category with translation $(\mathcal{D},T)$, and let
	$\epsilon,\zeta,\eta\in\Bbbk^\times$. If
	$\epsilon\zeta\eta=1$, then the original triangle is isomorphic to
	$X\xrightarrow{\epsilon f}Y\xrightarrow{\zeta g}Z
	\xrightarrow{\eta h}TX$, as shown by the commutative diagram
% segment-id: o014.li2.chapter4.triangulated-category-definition.g004
	\[\begin{tikzcd}
		X \arrow[r,"f"] \arrow[d,"1"'] &
		Y \arrow[r,"g"] \arrow[d,"\epsilon"] &
		Z \arrow[r,"h"] \arrow[d,"\epsilon\zeta"] &
		TX \arrow[d,"\epsilon\zeta\eta"] \\
		X \arrow[r,"\epsilon f"'] & Y \arrow[r,"\zeta g"'] &
		Z \arrow[r,"\eta h"'] & TX.
	\end{tikzcd}\]
\end{remark}

% segment-id: o014.li2.chapter4.triangulated-category-definition.p008
In the notation of Definition~\ref{def:morphism-with-degree}, a triangle can
also be written $X\to Y\to Z\xrightarrow{+1}$, or more pictorially as
% segment-id: o014.li2.chapter4.triangulated-category-definition.g005
$\begin{tikzcd}[row sep=small,column sep=tiny]
	& Y \arrow[rd] & \\
	X \arrow[ru] & & Z \arrow[ll,"{+1}"]
\end{tikzcd}$.
Thus a triangle may be pictured as an upward-spiraling chain, while a
morphism between triangles resembles the structure of deoxyribonucleic acid:
the arrows $\alpha$, $\beta$, $\gamma$, and so forth can be compared to the
hydrogen bonds between bases.
\index{deoxyribonucleic acid}

% segment-id: o014.li2.chapter4.triangulated-category-definition.d004
\begin{definition}[\mbox{Triangulated categories} and
\mbox{triangulated functors}]
\label{def:triangulated-cat}
	\index{triangulated category}
	\index{pretriangulated category}
	\index{functor!triangulated}
	\index{triangle!distinguished (exact)}
	A \emph{triangulated category} is an additive category with translation
	$(\mathcal{D},T)$ equipped with a class $\mathcal{H}$ of triangles, called
	\emph{distinguished triangles}, satisfying the following axioms.
% segment-id: o014.li2.chapter4.triangulated-category-definition.l002
	\begin{enumerate}[(TR1)]
		\setcounter{enumi}{-1}
% segment-id: o014.li2.chapter4.triangulated-category-definition.i004
		\item Every triangle isomorphic to a distinguished triangle is
		distinguished.
% segment-id: o014.li2.chapter4.triangulated-category-definition.i005
		\item For every $X\in\Obj(\mathcal{D})$, the triangle
		$X\xrightarrow{\identity_X}X\to0\to TX$ is distinguished.
% segment-id: o014.li2.chapter4.triangulated-category-definition.i006
		\item Every morphism $f:X\to Y$ can be extended to a distinguished
		triangle of the form $X\xrightarrow{f}Y\to Z\to TX$, where
		$Z\in\Obj(\mathcal{D})$.
% segment-id: o014.li2.chapter4.triangulated-category-definition.i007
		\item A triangle
		$X\xrightarrow{f}Y\xrightarrow{g}Z\xrightarrow{h}TX$ is distinguished
		if and only if its “counterclockwise rotation”
% segment-id: o014.li2.chapter4.triangulated-category-definition.q002
		\[ Y\xrightarrow{g}Z\xrightarrow{h}TX\xrightarrow{-Tf}TY \]
		is distinguished. The sign pattern may equivalently be chosen as
		$(-++)$, $(+-+)$, or $(---)$; see
		Remark~\ref{rem:triangle-signs}.
% segment-id: o014.li2.chapter4.triangulated-category-definition.i008
		\item Suppose the solid part of the following diagram is commutative:
% segment-id: o014.li2.chapter4.triangulated-category-definition.g006
		\[\begin{tikzcd}
			X \arrow[r] \arrow[d,"\alpha"'] &
			Y \arrow[r] \arrow[d,"\beta"] &
			Z \arrow[r] \arrow[dashed,d,"\gamma"] &
			TX \arrow[dashed,d,"T\alpha"] \\
			X' \arrow[r] & Y' \arrow[r] & Z' \arrow[r] & TX'.
		\end{tikzcd}\]
		If both rows are distinguished triangles, there exists a morphism
		$\gamma:Z\to Z'$ as indicated by the dashed arrow such that the whole
		diagram containing $\alpha$, $\beta$, $\gamma$, and $T\alpha$
		commutes. Thus one obtains a morphism of triangles.
% segment-id: o014.li2.chapter4.triangulated-category-definition.i009
		\item Suppose that distinguished triangles
% segment-id: o014.li2.chapter4.triangulated-category-definition.q003
		\begin{gather*}
			X\xrightarrow{f}Y\xrightarrow{h}Z'\to TX,\\
			Y\xrightarrow{g}Z\xrightarrow{k}X'\to TY,\\
			X\xrightarrow{gf}Z\xrightarrow{m}Y'\to TX
		\end{gather*}
		are given. There exists a distinguished triangle
% segment-id: o014.li2.chapter4.triangulated-category-definition.q004
		\[ Z'\xrightarrow{u}Y'\xrightarrow{v}X'\xrightarrow{w}TZ' \]
		such that the following diagram commutes:
% segment-id: o014.li2.chapter4.triangulated-category-definition.g007
		\[\begin{tikzcd}
			X \arrow[r,"f"] \arrow[equal,d] &
			Y \arrow[r,"h"] \arrow[d,"g"] &
			Z' \arrow[r] \arrow[d,"u"] & TX \arrow[equal,d] \\
			X \arrow[r,"gf"] \arrow[d,"f"'] &
			Z \arrow[r,"m"] \arrow[equal,d] &
			Y' \arrow[r] \arrow[d,"v"] & TX \arrow[d,"Tf"] \\
			Y \arrow[r,"g"] \arrow[d,"h"'] &
			Z \arrow[r,"k"] \arrow[d,"m"] &
			X' \arrow[r] \arrow[equal,d] & TY \arrow[d,"Th"] \\
			Z' \arrow[r,"u"'] & Y' \arrow[r,"v"'] &
			X' \arrow[r,"w"'] & TZ'.
		\end{tikzcd}\]
	\end{enumerate}

% segment-id: o014.li2.chapter4.triangulated-category-definition.p009
	Data $(\mathcal{D},T,\mathcal{H})$ satisfying only (TR0)--(TR4) are
	called a \emph{pretriangulated category}. The data $T$ and $\mathcal{H}$
	are often omitted from the notation.

% segment-id: o014.li2.chapter4.triangulated-category-definition.p010
	A \emph{triangulated functor} from $\mathcal{D}$ to $\mathcal{D}'$ is a
	functor $F:\mathcal{D}\to\mathcal{D}'$ between additive categories with
	translation that maps distinguished triangles to distinguished triangles.
	Morphisms between triangulated functors are the same as the morphisms in
	Definition~\ref{def:category-with-translation}; this also defines a
	quasi-inverse of a triangulated functor. If two triangulated functors
% segment-id: o014.li2.chapter4.triangulated-category-definition.g008
	$\begin{tikzcd}
		\mathcal{D} \arrow[shift left,r,"F"] &
		\mathcal{D}' \arrow[shift left,l,"G"]
	\end{tikzcd}$
	are quasi-inverse to one another, they are said to form an
	\emph{equivalence} between $\mathcal{D}$ and $\mathcal{D}'$.
\end{definition}

% segment-id: o014.li2.chapter4.triangulated-category-definition.p011
If $(\mathcal{D},T)$ is $\Bbbk$-linear, all the definitions above have
corresponding extensions.

% segment-id: o014.li2.chapter4.triangulated-category-definition.r002
\begin{remark}\index{octahedral axiom}
	In the notation of Definition~\ref{def:morphism-with-degree}, axiom (TR5)
	can be rewritten as the diagram
% segment-id: o014.li2.chapter4.triangulated-category-definition.g009
	\[\begin{tikzcd}[column sep=large,row sep=large]
		& Y' \arrow[dashed,rd,"v"] \\
		Z' \arrow[dashed,ru,"u"] \arrow[d,"+1" description] &&
		X' \arrow[dashed,ll,"+1" description,"w" inner sep=0.6em]
		\arrow[ldd,"+1" description,near end] \\
		X \arrow[rd,"f"'] && Z \arrow[u,"k"'] \\
		& Y \arrow[ru,"g"'] \arrow[luu,"h"',near start] &
		\arrow[from=1-2,to=3-1,crossing over,"+1" description,near start]
		\arrow[from=3-1,to=3-3,crossing over]
		\arrow[from=3-3,to=1-2,crossing over,"m",near end]
	\end{tikzcd}\]
	The dashed part represents the morphisms whose existence is asserted by
	the axiom. Four faces of this octahedron are distinguished triangles, and
	the other four faces are commutative; the details are left to the reader.
	For this reason, (TR5) is also called the \emph{octahedral axiom}.
\end{remark}

% [Yekutieli, Remark 5.6.12]
% segment-id: o014.li2.chapter4.triangulated-category-definition.r003
\begin{remark}[Duality]\label{rem:triangulated-op}
	Let $(\mathcal{D},T,\mathcal{H})$ be a pretriangulated (or triangulated)
	category. Regard $T^{-1}$ as an endofunctor of
	$\mathcal{D}^{\opp}$ and denote it by $S$. Define a class of triangles
	$\mathcal{K}$ on $(\mathcal{D}^{\opp},S)$ as follows. If
% segment-id: o014.li2.chapter4.triangulated-category-definition.q005
	\[ [X\xrightarrow{f}Y\xrightarrow{g}Z\xrightarrow{h}TX]
		\;\in\mathcal{H}, \]
	then its clockwise rotation
	$[T^{-1}Z\xrightarrow{-T^{-1}h}X\xrightarrow{f}Y\xrightarrow{g}Z]
	\in\mathcal{H}$ may be regarded as the following triangle in
	$\mathcal{D}^{\opp}$:
% segment-id: o014.li2.chapter4.triangulated-category-definition.q006
	\begin{equation}\label{eqn:op-triangle}
		Z\xrightarrow{g^{\opp}}Y\xrightarrow{f^{\opp}}X
		\xrightarrow{-(T^{-1}h)^{\opp}}SZ.
	\end{equation}
	Define $\mathcal{K}$ to be the class of all triangles of the form
	\eqref{eqn:op-triangle}. It is straightforward to verify that
	$(\mathcal{D}^{\opp},S,\mathcal{K})$ is a pretriangulated (or
	triangulated) category. It is also clear that applying this procedure
	twice recovers the original $(\mathcal{D},T,\mathcal{H})$.
\end{remark}

% segment-id: o014.li2.chapter4.triangulated-category-definition.p012
The notions of pretriangulated subcategory and triangulated subcategory will
be introduced in
Definition--Proposition~\sourcecrossref{def:triangulated-subcat}{in \S~4.2};
we postpone their discussion until then.
